% Transcription — fonds Alexandre Grothendieck, shelfmark 97, batch 1
% (pages 1-18, the whole folder). Established from the facsimile
% archives/batches/97/batch-01.pdf (a mirror of
% https://grothendieck.umontpellier.fr/97.pdf), per the skill
% transcribe-grothendieck, read from the tiles under archives/tiles/97/.
% The batch file carries no cover sheet: PDF page and archivists' page
% coincide.
%
% Fourteen of the eighteen pages are transcribed. Four are not: page 2, a
% list of mathematicians' names in his hand with brackets and boxes (a
% distribution list, no mathematics); page 3, a sheet of IHES letterhead
% carrying only the abandoned opening of a letter to Atiyah; page 5, blank;
% page 14, the blank verso of page 13, showing only its bleed-through. The
% gaps in the numbering are the record.
%
% What the folder holds, in the archivists' order.
%
% Page 1 is the letter of the catalogue title: J.-P. Serre to Grothendieck,
% on Collège de France letterhead, « Paris, le 4 Nov. 66 », in red ink —
% Serre has not looked at the problem, but sends a « Notice » proving the
% same thing, which cites a paper of Milnor on flat SL_2 bundles over a
% surface. It is a letter by someone other than Grothendieck, and it is
% handled as folder 134-1 handles Deligne's covering note (transcribed, with
% a note saying whose it is) and as folder 29 handles the Murre letters
% (transcribed whole): it is short, it is addressed to him, and its one
% mathematical sentence names the question the manuscript pages 6-12 work
% on. It is given whole. The letterhead is institutional; no private postal
% address or telephone number appears on it, so nothing is withheld.
%
% Page 4, blue ink, carries his own title « Classe de Chern de la
% représentation régulière et ses multiples »: for G cyclic of order p the
% total Chern class of the regular representation is the product of
% (1 + λξ_0), whence c_{p-1}(r_G) = ξ_0^{p-1} and, by restriction to a
% subgroup of order p, the non-vanishing of H^{Nk}(G)(p) for every N with
% k = 2n'(p-1). The page computes the product over F_p^* as 1; by Wilson's
% theorem it is -1. Flagged in a note, not repaired.
%
% Pages 6, 8, 10, 12 are four manuscript sheets in black ink, numbered by
% him ①, (none), ②, ③, written on the backs of the typescript sheets that
% follow them. The subject is the one of Serre's letter: X = D/Γ a compact
% surface of genus g >= 2 uniformised by the disc, G = SL(2,R)/{±1} acting
% on D with stabiliser S, the universal covering S~ -> S and the lift Γ~ of
% Γ to SL(2,R); the tangent bundle T_X and a real bundle Σ_0 of rank 2 with
% Σ_0 ⊗ Σ_0 related to T_X; the class γ_1 ∈ H^2(B_G, Z), γ_1 = 2c_1 on B_S,
% and its pull-back γ_1^X; the computation γ_1^X = 2 - 2g = EP(X), checked
% again by passing to a finite étale covering of degree d; page 12 reads
% γ_1 as the class of the extension of Γ by Z given by the universal
% covering of G, looks at H^2(Γ, Z/2Z), and in a « Remarque (4) » notes that
% on X the real bundle built from T_X plus two trivial summands carries an
% integrable connection. The hand on pages 8 and 12 is the fastest of the
% folder and much of their connective prose is left \ill{}; the displayed
% formulas and the diagrams of pages 6 and 10 read.
%
% Pages 7, 9, 11, 13 are four sheets of a French typescript on the Brauer
% group of a fibration X -> Y over a curve — typed pages 23, 17, 19, 20, so
% out of order — with corrections in his hand: Greek letters, tildes and
% Picard schemes inked in, interlinear insertions (« ainsi que (4.16),
% puisque k(y) est de dimension <= 1 ... », « [24] H^1(k,G)=0 »), a
% handwritten line of the exact-sequence argument and a large block on
% page 11 cancelled by diagonal strokes. The folder 29 rule applies: a typed
% verso is kept when it carries something in his hand. They are given in the
% archivists' order: page 9 (typed 17) defines B = i_*(i^*P), the skyscraper
% sheaves E and F and their geometric fibres; page 11 (typed 19) the
% Tate-Shafarevich group Ш(Y,B) and the sequence (4.17); page 13 (typed 20)
% the multiplicities μ, ν and the exact sequence (4.21) killing H^1(y,E_y)
% by d_y = pgcd(a ν); page 7 (typed 23) the two special cases, Y a curve
% over an algebraically closed field, Y a proper curve over a finite field.
% The typescript's Ш is typed as an overstruck III and is set \mathrm{III};
% its underlined sheaf letters are set \underline{} inside mathematics.
% Nothing on these sheets connects them to Chern classes; they are the paper
% he wrote on.
%
% Page 15, blue ink on the verso of page 16: H^*(B_G, R) for a compact
% connected Lie group, computed by the bicomplex of invariant forms on the
% G^{p+1}/G, degeneration of the spectral sequence, transgressive elements
% T and the theorem Sym^*(T) ≅ H^*(B_G, R).
%
% Page 16 is a typescript sheet of a letter — no heading, no signature on
% this leaf, the addressee called « tu » and a letter « du 26.3 » mentioned
% — on the categories C(X,A) as homotopy invariants of X (Hurewicz,
% Artin-Mazur), cancelled by three long diagonal strokes and carrying one
% insertion in his hand (« squelettes, images des groupes d'homotopie
% supérieurs »). Kept by the same rule; no author or date is supplied.
%
% Page 17, blue ink: for a connected (complex reductive, by an insertion
% whose reading is doubtful) Lie group and a prime p, the conditions « G
% sans p-torsion », « B_G sans p-torsion », the surjectivity of
% Sym(B^ ⊗ F_p) -> H^*(G/B, F_p), their bis variants, and the web of
% implications between them (Hopf, Borel transgression, Leray), drawn as a
% small diagram at the foot.
%
% Page 18 is a sheet of the EGA IV typescript, page IV-1102, (21.6.5) to
% (21.6.6) on the cycle of a positive divisor, with corrections in his hand
% and the marginal notes « devrait figurer dans 21.2 » and « réf ! ». As in
% folder 29 page 105, insertions (typed over or handwritten) are \add{} and
% deletions \struck{}; typed underlinings are \emph{}.
%
% Main doubtful readings: the group name on page 6 (« GP(2,R) », read as
% written although it is Sl(2,R)/{±1}); the whole of page 8 below its first
% formulas; page 12 almost throughout; the insertion « complexe réductif »
% on page 17; the marginal « (sans noeth.) » on page 18; Hom in
% H^2(G,Z) ≅ Hom(G,Q/Z) on page 4.
%
% The dating is the inventory's own, for the folder. Serre's letter is
% dated on the page, 4 November 1966.
%
% Pass: Opus 5 (claude-opus-5), 2026-09-14 — first pass, unchecked against the pages by a human.
\documentclass[11pt,a4paper]{article}
% Le préambule commun à toutes les transcriptions du fonds Grothendieck.
%
% Il tient en une idée : **l'incertitude se compose, elle ne se résout pas.**
% Une transcription qui lisse un mot douteux a détruit la seule chose pour
% laquelle on la faisait — un lecteur doit pouvoir savoir, à chaque endroit,
% ce qui était sur le feuillet et ce qui a été ajouté. Les six macros qui
% suivent sont donc l'appareil critique, et non de la décoration.
%
% Les mêmes commandes sont reconnues par `scripts/render.mjs`, qui produit la
% vue de lecture du site. Une macro ajoutée ici doit l'être aussi là-bas,
% faute de quoi le rendu échouera bruyamment — ce qui est voulu : un rendu
% silencieusement faux est pire qu'un rendu absent.

\ProvidesPackage{grothendieck}[2026/08/08 Transcriptions du fonds Grothendieck]

\usepackage{amsmath,amssymb,amsthm}
\usepackage{mathrsfs}
\usepackage{stmaryrd}
% Les diagrammes commutatifs sont partout dans ce fonds : c'est la langue
% même de la géométrie algébrique de Grothendieck, et une transcription qui
% les rendrait en prose ne transcrirait rien.
\usepackage{tikz-cd}
% Français par défaut — c'est la langue des feuillets — mais l'anglais est
% chargé aussi : la traduction et le résumé sont des documents anglais, et la
% typographie française (espaces avant les deux-points) y serait une faute.
% Ces documents basculent par \selectlanguage{english} en tête de corps.
\usepackage[english,french]{babel}
\usepackage{fontspec}
\usepackage{geometry}
\geometry{a4paper,margin=25mm,marginparwidth=28mm}
\usepackage{marginnote}
\usepackage{xcolor}
% ulem, not soul. `soul`'s \st and \ul are built on a character-by-character
% scan that XeLaTeX's font machinery breaks: the document fails with an
% undefined control sequence rather than degrading. ulem does the same two jobs
% and is Unicode-engine safe. `normalem` stops it hijacking \emph, which the
% transcriptions use for Grothendieck's own emphasis.
\usepackage[normalem]{ulem}
\usepackage{hyperref}
\hypersetup{colorlinks=true,linkcolor=black,urlcolor=black,pdfborder={0 0 0}}

% --- Métadonnées du lot -----------------------------------------------------
% Reprises telles quelles par le script de rendu, qui les affiche en tête de
% la vue de lecture. Elles fixent ce que le fichier prétend couvrir : sans
% elles, une transcription flotte sans point d'attache dans un fonds de seize
% mille feuillets.

\newcommand{\folder}[1]{\gdef\grothFolder{#1}}
\newcommand{\batch}[1]{\gdef\grothBatch{#1}}
\newcommand{\leaves}[2]{\gdef\grothFirst{#1}\gdef\grothLast{#2}}
\newcommand{\foldertitle}[1]{\gdef\grothFoldertitle{#1}}
\gdef\grothFolder{?}\gdef\grothBatch{?}\gdef\grothFirst{?}\gdef\grothLast{?}\gdef\grothFoldertitle{}

% --- Le résumé --------------------------------------------------------------
%
% Il ouvre la lecture modernisée, sous le titre « Résumé », et il est composé
% en petit. Ce n'est pas un ornement : c'est le seul passage du document qui
% ne soit pas des mathématiques, et un lecteur doit pouvoir le reconnaître
% avant de l'avoir lu — puis le sauter s'il connaît déjà le sujet, ou s'y
% arrêter s'il ne le connaît pas. Le filet qui le suit marque la reprise du
% corps.

\newenvironment{resume}
  {\small\setlength{\parskip}{0.45em}}
  {\par\medskip\hrule\medskip}

% --- Mots-clés ---------------------------------------------------------------
%
% En anglais, en clôture du résumé de la lecture modernisée : le vocabulaire
% moderne sous lequel chercher ce que les feuillets construisent. Le manifeste
% du site (`npm run manifest`) les extrait pour étiqueter la cote entière —
% cette ligne est la source unique des tags, il n'y a pas de fichier à côté.
\newcommand{\keywords}[1]{\par\smallskip\noindent
  {\footnotesize\textsc{Keywords} --- \textit{#1}}\par}

% --- L'appareil critique ----------------------------------------------------

% Le numéro de feuillet, en marge. C'est l'ancre qui permet de régler un
% désaccord sur une formule en regardant *un* feuillet plutôt qu'une chemise
% de six cents. Le site s'en sert aussi pour tourner les pages du fac-similé
% quand on fait défiler la transcription.
\newcommand{\leaf}[1]{%
  \marginnote{\footnotesize\color{gray!70}\textsf{#1}}%
  \label{leaf:#1}\ignorespaces}

% Illisible : on ne devine pas. Un mot inventé qui se lit comme les autres est
% le pire résultat possible.
\newcommand{\ill}{\textcolor{red!60!black}{[\,\ldots\,]}}

% Lecture douteuse : le mot est donné, et signalé comme incertain.
\newcommand{\uncertain}[1]{\uline{#1}}

% Ajout de l'éditeur — une lettre manquante, un mot sous-entendu.
\newcommand{\add}[1]{\textcolor{blue!50!black}{[#1]}}

% Biffé par Grothendieck lui-même. Ce qu'il a rayé fait partie du document :
% une rature dit souvent où la pensée a changé de direction.
\newcommand{\struck}[1]{\sout{#1}}

% Note du transcripteur, sur l'état matériel du feuillet.
\newcommand{\note}[1]{\par\smallskip{\small\color{gray!60!black}\textit{[#1]}}\par\smallskip}

% Note marginale de Grothendieck, distincte du corps du texte.
\newcommand{\marginal}[1]{\par\smallskip\noindent\hspace{1em}%
  {\small\color{gray!60!black}#1}\par\smallskip}

% --- Titre ------------------------------------------------------------------

\AtBeginDocument{%
  \begin{center}
    {\large\bfseries Fonds Alexandre Grothendieck --- cote n\textsuperscript{o}~\grothFolder}\\[2pt]
    {\itshape\grothFoldertitle}\\[4pt]
    {\small feuillets \grothFirst--\grothLast\ (lot \grothBatch)}\\[2pt]
    {\footnotesize\color{gray!60!black}Transcription établie d'après le fac-similé numérisé
     par l'Université de Montpellier.\\
     Les lectures incertaines sont soulignées, les illisibles notées [\,\ldots\,],
     les ajouts entre crochets.}
  \end{center}
  \vspace{1em}}

\endinput


\folder{97}
\batch{1}
\pages{1}{18}
\dating{[vers 1966]}
\watermark{Édition de démonstration}
\foldertitle{Classes de Chern des représentations : lettre (1966), notes manuscrites (s.d.)}

\begin{document}

\section*{J.-P.\ Serre à A.\ Grothendieck (page 1)}

\page{1}
\note{Lettre de Serre, à l'encre rouge, sur papier à en-tête du Collège de
France, chaire d'algèbre et géométrie. Elle n'est pas de Grothendieck ; elle
est donnée entière, comme la note de Deligne du dossier 134-1 et les lettres
de Murre du dossier 29, parce qu'elle nomme la question des feuillets 6 à 12}

Paris, le 4 Nov.\ 66

Cher Grothendieck,

Je n'ai pas regardé ton pb. Mais je suis tombé sur une « Notice » que je
t'envoie ci-joint, démontrant la \uncertain{même chose} que toi.

La notice renvoie à un papier de Milnor, que j'ai, où celui-ci regarde
justement les fibrés à groupe $\mathrm{Sl}_2$ qui sont plats, et de base une
surface. C'est très joli.

Bien à toi

J-P.\ Serre

\section*{Classe de Chern de la représentation régulière et ses multiples (page 4)}

\note{titre de sa main, en tête du feuillet}

\page{4}
$G$ groupe cyclique d'ordre $n$,
\[
H^{*}(G, \mathbb{Z})^{+} \simeq
\mathrm{Sym}_{\mathbb{Z}/n\mathbb{Z}}(\check{G})^{+} \simeq
\mathbb{Z}/n\mathbb{Z}[\xi]^{+}
\]
\note{le $\mathbb{Z}$ de $H^{*}(G,\mathbb{Z})$ est surchargé}
\marginal{au-dessus de $\check{G}$ : dual sur $\mathbb{Z}/n\mathbb{Z}$}
où $\xi$ est un générateur de $H^{2}(G,\mathbb{Z}) \simeq
\mathrm{Hom}(G,\mathbb{Q}/\mathbb{Z}) \simeq \check{G}$.
\note{« $\mathrm{Hom}$ » est une lecture incertaine : le signe ressemble à un $H$ surmonté d'une étoile}

La représentation régulière \add{sur $\mathbb{C}$} \struck{est} correspond à
la somme des représentations de degré 1 \uncertain{associées} aux
$\xi \in \check{G}$. On a donc
\[
c(r_{G}) = \prod_{\xi \in \check{G}} (1 + \xi)
= \prod_{\lambda \in \mathbb{Z}/n\mathbb{Z}} (1 + \lambda \xi_{0}) .
\]

Supposons $n = p$ premier. On trouve
\[
\begin{cases}
c_{p}(r_{G}) = 0 \\
c_{p-1}(r_{G}) = \prod_{\xi \in \check{G} - \{0\}} \xi
= \Bigl(\prod_{\lambda \in \mathbb{F}_{p}^{*}} \lambda\Bigr) \xi_{0}^{p-1}
= \xi_{0}^{p-1}
\end{cases}
\]
(indépendant du choix de $\xi_{0} \in \check{G}$ !)
\note{par le théorème de Wilson le produit des $\lambda \in
\mathbb{F}_{p}^{*}$ vaut $-1$, et non $1$ ; la page écrit $\xi_{0}^{p-1}$ et
la suite en dépend, au signe près seulement}

D'où
\[
\begin{cases}
c_{i}(N r_{G}) = 0 & \text{si } i > N(p-1) \\
c_{N(p-1)}(N r_{G}) = \xi^{N(p-1)} & (\xi \in \check{G})
\end{cases}
\]

Soit maintenant $G$ un groupe fini, d'ordre $n$, divisible par $p$,
\struck{Alors \ill{} tous} et soit $H$ un sous-groupe d'ordre $p$.
Considérons $n' = n/p$, alors $r_{G}|H = n' r_{H}$ donc $(N r_{G})|H = (N n')
r_{H}$, d'où \struck{$(N r_{G})|H$}
\[
(c_{i}(\uncertain{n}\, r_{G}))|H = c_{i}(N n' r_{H}) ,
\]
et par suite $c_{N n'(p-1)}(r_{G})|H \neq 0$, a fortiori $c_{N n'(p-1)}(r_{G})
\neq 0$.
\note{dans ces deux dernières lignes la page écrit $r_{G}$, sans le facteur
$N$ des lignes précédentes, et un $n$ minuscule qui est vraisemblablement le
$N$}

Donc
\[
\boxed{H^{Nk}(G)(p) \neq 0 \text{ pour tout entier } N, \text{ où } k = 2n'(p-1)}
\]

\section*{Surfaces de genre $g$ et la classe $\gamma_{1}$ (pages 6, 8, 10, 12)}

\note{quatre feuillets à l'encre noire, numérotés par lui \emph{1}, (sans
numéro), \emph{2}, \emph{3}, écrits au dos des feuillets du tapuscrit qui
les suivent dans le dossier}

\page{6}

\begin{tikzcd}[column sep=small, row sep=small]
\widetilde{S} \arrow[rr] & & S \\
\widetilde{G} = \mathrm{Sl}(2,\mathbb{R}) \arrow[rr] \arrow[d, no head] & &
G = \mathrm{GP}(2,\mathbb{R}) = \widetilde{G}/(+1,-1)
\arrow[d, no head] \\
\widetilde{\Gamma} \arrow[rr, "\sim"] & & \Gamma
\end{tikzcd}

\begin{tikzcd}
V \\
D \arrow[d] \\
X
\end{tikzcd}

\note{« $\mathrm{GP}(2,\mathbb{R})$ » est lu tel qu'écrit, lecture incertaine. Les traits verticaux sous $\widetilde{G}$ et $G$ n'ont pas de pointe
lisible ; ils portent vraisemblablement les inclusions $\widetilde{\Gamma}
\subset \widetilde{G}$, $\Gamma \subset G$. Sous $\widetilde{G}/(\ldots)$ un
mot biffé, illisible, avant $(+1,-1)$}

$T_{X}$, $T_{D}$ avec op.\ de $\Gamma$ \struck{\ill{}} induit par $T_{D}$
avec op.\ de $G$, déduit par op.\ de $S = G_{x}$ sur $T_{D}(x)$.

$T_{D}$ induit par $T_{D}$ avec op.\ de $\widetilde{G}$, déduit par op.\ de
$\widetilde{S} = \widetilde{G}_{x} \to S$ sur $T_{D}(x)$.

$E$, défini par rep.\ de $\widetilde{\Gamma}$ induite par $\widetilde{G}$
op.\ sur $V$, déduite par fibré vect.\ $D \times V$ avec op.\ diagonale de
$\widetilde{\Gamma}$ : $\gamma \cdot (x,v) = (\gamma x, \gamma v)$, induite
par op.\ diagonale de $\widetilde{G}$ sur $D \times V$, déduite par $V$ vue
comme \ill{} rep.\ de $\widetilde{S}$ sur $V$.

On a
\[
\boxed{X = X_{0}^{(2)}}
\qquad \text{caractère de } X_{0} : e_{1} + e_{-1} ,
\qquad \text{caractère de } X : \struck{\ill{}}\ e_{2} + e_{-2} ,
\qquad X = X_{0}^{2} - 2 .
\]

N.B.\ Pour toute \struck{\ill{}} représentation linéaire réelle
\struck{rep.\ complexe de \ill{}} de $\widetilde{S}$, i.e.\
$\widetilde{\Gamma} \simeq \Gamma \subset \widetilde{G}$, conduisant à un
fibré vectoriel réel \struck{triv.} sur $X$ \add{(rep.\ complexe --- mais pas
nécessairement \ill{})}. Cette correspondance \ill{}
\marginal{dans la marge, à l'encre brune, un mot biffé, \uncertain{Lusin}}

\page{8}
\note{suite du feuillet 6, d'une écriture très rapide ; une grande partie
de la prose est perdue}

\uncertain{respecte} les opérations tensorielles, l'opération de
\uncertain{conjugaison} \ill{} ; où $X_{0}$ correspond à un fibré
\uncertain{vectoriel complexe} de rang 1, et $X$ correspond au
\uncertain{carré} \ill{}.

\struck{Par suite} Si $\Sigma = T_{X}$, $\Sigma_{0}$ sont les fibrés
vectoriels \add{complexes} \struck{réels} inversibles, leurs
\struck{classes de Chern} \add{\ill{}} $\widetilde{c}_{1}$, $c_{1}$ sont
reliées par
\[
\widetilde{c}_{1} = 2 c_{1} .
\]

D'autre part on connaît aussi les classes \struck{de Chern}
\add{caractéristiques} des deux \uncertain{fibrés envisagés}. Comme
\add{\ill{}} $\Sigma \simeq T_{X}$, \ill{}
\[
\gamma_{1} = EP(X) = 2(1-g) , \quad \ill{}
\]
\[
c_{1} = c_{1}^{\mathbb{R}}(\Sigma_{0}) = (1-g)
\]
\note{le $\widetilde{c}_{1}$ de la ligne précédente est ici écrit comme un
$\gamma_{1}$ ; c'est la notation du feuillet 10}

N.B.\ Le fibré vectoriel \add{réel sur $X$} $\Sigma_{0}$ \struck{\ill{} est}
\ill{}, d'une \ill{}, par passage au quotient d'un \ill{} vectoriel sur $D$,
à l'aide d'une \uncertain{représentation réelle} (pas complexe !) de
$\Gamma$, \ill{}
\note{la fin du feuillet, sept lignes, ne se lit que par mots isolés : « sur
$D$ », « de $\Gamma$ », « de $\widetilde{S}$ », « structure », « les fibrés
$\Sigma_{0}$ », « $T_{X}$ », « d'une connexion intégrable »}

\marginal{Re.\ NB. \ill{} $\Sigma_{0} \otimes_{\mathbb{R}} \Sigma_{0} \simeq
T_{X} + 2 \cdot 1_{\mathbb{R}}$ ; comme $T_{X} + 2 \cdot 1_{\mathbb{R}}$
admet une connexion intégrable \ill{}}
\note{note écrite verticalement dans la marge de gauche ; lecture très
incertaine au-delà des formules}

\page{10}
\[
\Gamma \subset G \quad \text{d'où :}
\]

\begin{tikzcd}[column sep=small, row sep=small, nodes={font=\scriptsize}]
& H^{*}(B_{G},\mathbb{Z}) \arrow[rr] \arrow[dl] \arrow[d, "\wr"] & &
H^{*}(B_{\Gamma},\mathbb{Z}) \simeq H^{*}(X,\mathbb{Z}) \\
H^{*}(B_{\widetilde{G}},\mathbb{Z}) \arrow[d, "\wr"] &
H^{*}(B_{S},\mathbb{Z}) \simeq \mathbb{Z}[\gamma_{1}] \arrow[dl] & & \gamma_{1}^{X} \\
H^{*}(B_{\widetilde{S}},\mathbb{Z}) \simeq \mathbb{Z}[c_{1}] & & &
\end{tikzcd}

\note{une flèche ondulée mène de $\mathbb{Z}[\gamma_{1}]$ à
$\gamma_{1}^{X}$, écrit sous $H^{*}(X,\mathbb{Z})$}

\uncertain{Car} $\gamma_{1} = 2c_{1}$

d'où
\[
\boxed{\gamma_{1}^{X} \in H^{2}(X,\mathbb{Z}) \simeq \mathbb{Z}}
\]

À calculer $\gamma_{1}^{X}$ ??

\struck{Sat}

\uncertain{Cependant} \ill{} $\widetilde{\Gamma}$ \ill{}

\begin{tikzcd}
B_{\Gamma} = X \arrow[r] \arrow[dr, dashed] & B_{G} \\
& B_{\widetilde{G}} \arrow[u]
\end{tikzcd}

\ill{} $=$
\[
\boxed{\tfrac{1}{2}\gamma_{1}^{X} = c_{1}^{X} \in H^{2}(X,\mathbb{Z}) = \mathbb{Z}}
\]
dont on \ill{} \add{\ill{} classe de Chern $c_{1}^{\mathbb{R}}$ d'un} \ill{}
\uncertain{fibré vectoriel réel} $\Sigma$ sur $B_{\widetilde{G}}$, de rang 2.

Le calcul donne $c_{1} = 1 - g$, donc
\[
\boxed{\gamma_{1}^{X} = 2(1-g) = 2 - 2g}
\qquad \{= EP(X) \neq 0
\]

Dans le cas \ill{} qu'il existe \ill{} $\Gamma'$ de $\Gamma$ \ill{},
d'indice $d$ :

\begin{tikzcd}
B_{\Gamma'} = X' \arrow[r, "p"] & B_{\Gamma} = X \arrow[r] & B_{G}
\end{tikzcd}

($p$ : rev.\ étale de degré $d$). On a aussi $\gamma_{1}^{X'} =
p^{*}(\gamma_{1}^{X})$, donc \ill{} $\gamma_{1}^{X'} = d\,\gamma_{1}^{X}$.
Mais aussi, comme bien connu, $EP(X') = d\,EP(X)$ ; et comme
$\gamma_{1}^{X'} = EP(X')$, on conclut $\gamma_{1}^{X} = EP(X)$. OK.
\note{la lecture de ce dernier paragraphe suit les formules, qui se lisent ;
les mots qui les relient sont en partie devinés, en particulier « comme
$\gamma_{1}^{X'} = EP(X')$ »}

\page{12}
\note{écriture la plus rapide du dossier ; seuls les formules et quelques
mots se lisent}

N.B.\ Le rev.\ universel de $G$ est de \struck{\ill{}} \add{noyau}
$\mathbb{Z}$. \uncertain{Je} \ill{} donne une extension \add{\ill{}} de
$\Gamma$ par $\mathbb{Z}$, \ill{}
\[
H^{2}(\Gamma,\mathbb{Z}) \simeq H^{2}(X,\mathbb{Z}) .
\]
\ill{} que c'est $\gamma_{1}$. Plus généralement, \ill{} $G$ \ill{} $H$),
\ill{} $G$, $H$ \ill{}), \ill{}
\[
H^{2}(B_{G}, H) \quad \ill{}
\]
Dans \ill{} $S =$ \ill{}, \ill{} $\widetilde{S} =$ rev.\ univer\add{sel}
\ill{} $\widetilde{S}$, de \ill{} $\mathbb{Z}$, \ill{}
\[
B_{H} \to B_{E} \to B_{G} \qquad \ill{}\ H^{2}(B_{S},\mathbb{Z})\ \ill{}\
\text{Chern}.
\]
\ill{} l'inter\ill{} \ill{}
\[
\gamma_{1}^{X} = \struck{\ill{}}\ \ill{}\ H^{2}(X,\mathbb{Z}) \simeq
H^{2}(\Gamma,\mathbb{Z}) .
\]
\[
H^{2}(\Gamma, \mathbb{Z}/2\mathbb{Z}) = \mathbb{Z}/2\mathbb{Z}
\quad \struck{\ill{}}
\]
\ill{} $\mathbb{Z}/2\mathbb{Z}$ \ill{} $\widetilde{G} \to G$ \ill{} !
Comme $EP(X) = 2(1-g)$ \ill{} $\widetilde{F}$ ! \ill{} !

(4) \uncertain{Remarque}, \ill{}, à la situation (1), \ill{} $g \geq 2$,
fibré tangent $T_{X}$, le fibré vectoriel \struck{complexe} \add{réel de
\ill{}} $(T_{X} \struck{\ill{}}) + 2 \cdot 1_{\mathbb{R}}$ sur $X$ est muni
d'une connexion intégrable, \ill{}

\section*{Tapuscrit sur le groupe de Brauer d'une fibration (pages 7, 9, 11, 13)}

\note{feuillets d'un tapuscrit, paginés 23, 17, 19 et 20 à la machine ou à
la main en haut à droite, donnés ici dans l'ordre des archivistes. Ils
portent des corrections de sa main : les insertions sont données par
\add{}, les mots tapés ou biffés par \struck{}, les soulignements du
tapuscrit par l'italique. Le Ш de Tate--Chafarévitch est tapé comme un III
barré et rendu $\mathrm{III}$ ; les lettres grecques et les tildes sont
presque tous ajoutés à la main}

\page{7}
\note{p.\ 23 du tapuscrit}

\[
0 \to \underline{\mathbb{Z}}/\delta''\underline{\mathbb{Z}} \to
\mathrm{III}'(Y,\underline{A}) \to \mathrm{III}'(Y,\underline{B}) \to 0 .
\tag{4.38}
\]

Explicitons deux cas particuliers importants :

\emph{Cas particulier} \struck{1} \add{(4.6) :} \emph{$Y$ est une courbe
algébrique sur un corps algébriquement clos.} Alors la \struck{relation
(4.13) est automatique la} plus grande partie de la discussion précédente
devient triviale. La relation (4.13) est automatiquement satisfaite, et on a
(4.12 bis) :
\[
\mathrm{Br}(X) \xrightarrow{\ \sim\ } H^{1}(Y,\underline{B}) ;
\tag{4.39}
\]
lorsque de plus les $\delta_{y}$ sont égaux à 1, i.e.\ \add{si} $X$ n'a pas
de fibre multiple sur $Y$, alors le deuxième membre s'explicite
\struck{comme étant égal à} par
\[
0 \to \underline{\mathbb{Z}}/\delta\underline{\mathbb{Z}} \to
H^{1}(Y,\underline{A}) \to H^{1}(Y,\underline{B}) \to 0 ,
\tag{4.40}
\]
compte tenu que \struck{dans} l'on a ici nécessairement $\delta = \delta' =
\delta''$.

\emph{Cas particulier} (4.7) : \emph{$Y$ est une courbe algébrique propre sur
un corps fini $k$.} Ici encore, (4.13) \struck{est} \add{\struck{\ill{}}}
automatiquement satisfaite, $k$ étant parfait, \add{ainsi que (4.16),
puisque $k(y)$ est de dimension $\leq 1$, d'où $H^{3}(k(y),\mathbb{G}_{m}) =
0$ ; de plus} \struck{et} on aura $\delta = \delta'$ \struck{\ill{}}, et de
même pour les invariants hensélisés et strictement hensélisés. Comme on a
encore $\mathrm{Br}(Y) = 0$, on trouve ici (4.4) :
\[
0 \to \mathrm{Br}(X) \to \mathrm{III}(Y,\underline{B}) \to
\underline{\mathbb{Z}}/\Delta\underline{\mathbb{Z}} \to 0 ,
\qquad \text{où } \Delta \mid \delta .
\tag{4.41}
\]
\struck{\ill{}} D'autre part, si les $\delta_{y}$ sont égaux à 1, i.e.\ si
le morphisme déduit de $f : X \to Y$ en passant à la clôture algébrique de
$k$ n'a pas de fibre multiple, alors on trouve la suite exacte
\[
0 \to \underline{\mathbb{Z}}/\delta''\underline{\mathbb{Z}} \to
\mathrm{III}(Y,\underline{A}) \to \mathrm{III}(Y,\underline{B}) \to 0 ,
\tag{4.42}
\]
ce qui est un cas particulier de (4.37), compte tenu de\struck{s} \add{la}
relation\struck{s} suivante\struck{s}
\[
\bar{\delta}_{y} = \delta_{y}
\tag{4.43}
\]
pour tout point fermé $y$ de $Y$, qui implique ici $\bar{\delta}_{y} = 1$
d'où $\mathrm{III}'(Y,\underline{A}) = \mathrm{III}(Y,\underline{A})$. La
démonstration de\struck{s} \add{l'} égalité\struck{s} (4.43), qui est une
conséquence facile du théorème de LANG \add{([24] $H^{1}(k,G) = 0$)}
pour le corps fini $k(y)$ et le schéma de Picard \struck{Con}connexe
\add{$G =$} $\underline{\mathrm{Pic}}^{0}_{X_{y}/y}$, est laissé\struck{e}
au lecteur comme un exercice plaisant et délectable.

\page{9}
\note{p.\ 17 du tapuscrit}

$H^{3}(Y,\underline{\mathbb{G}}_{m}) \simeq \underline{\mathbb{Q}}/
\underline{\mathbb{Z}}$, d'où le corollaire.

Soit \struck{\ill{}}
\[
i : \eta \longrightarrow Y
\]
l'inclusion du point générique, et posons
\[
\underline{B} = i_{*}(i^{*}(P)) ,
\tag{4.8}
\]
de sorte qu'on a un homomorphisme canonique
\[
P \longrightarrow \underline{B} ,
\tag{4.9}
\]
et on posera
\[
\underline{E} = \mathrm{Ker}(P \to \underline{B}) , \qquad
\underline{F} = \mathrm{Coker}(P \to \underline{B}) .
\tag{4.10}
\]
Les faisceaux $\underline{E}$ et $\underline{F}$ sont des « faisceaux
gratte-ciel », qui sont donc connus lorsqu'on connait leurs restrictions aux
points fermés de $Y$. Leurs fibres géométriques au dessus d'un tel point $y$
s'obtiennent comme les noyaux et conoyau de l'homomorphisme
\[
\add{\mathrm{Pic}(\widetilde{X}) =}\ \struck{\mathrm{Pic}}\
\mathrm{Pic}(\widetilde{X}/\widetilde{Y}) \longrightarrow
\mathrm{Pic}(\widetilde{X}_{\tilde{\eta}}/\tilde{\eta}) ,
\]
où les $\sim$ dénotent l'effet de la localisation stricte relativement à un
point géométrique $\bar{y}$ au dessus de $y$. \struck{\ill{}} Compte tenu que
le morphisme précédent se factorise par le sous-groupe
$\mathrm{Pic}(\widetilde{X}_{\eta})$, et que $\mathrm{Pic}(\widetilde{X})
\to \mathrm{Pic}(\widetilde{X}_{\tilde{\eta}})$ est surjectif puisque
$\widetilde{X}$ est régulier, on trouve
\begin{align}
\struck{\ill{}}\quad \underline{F}_{\bar{y}} &=
\mathrm{Pic}(\widetilde{X}_{\eta}/\tilde{\eta}) /
\mathrm{Pic}(\widetilde{X}_{\tilde{\eta}}) \tag{4.11} \\
\underline{E}_{\bar{y}} &= \mathrm{Ker}\bigl(\mathrm{Pic}(\widetilde{X})
\longrightarrow \mathrm{Pic}(\widetilde{X}_{\tilde{\eta}})\bigr) . \tag{4.12}
\end{align}
Compte tenu que $\mathrm{Br}(\tilde{\eta}) = 0$ si $k(y)$ est parfait, et que
c'est un groupe de $p(y)$-torsion en tous cas, où $p(y)$ est l'exposant
caractéristique de $k(y)$, on trouve que $\underline{F}_{\bar{y}}$ est un
groupe de $p(y)$-torsion, nul si $k(y)$ est parfait. Désignant par
$\delta_{y}$, $\delta'_{y}$, $\delta''_{y}$ les entiers analogues à
$\delta$, $\delta'$, $\delta''$ pour la situation localisée $\widetilde{X}
\to \widetilde{Y}$, un argument déjà \struck{\ill{}} fait plus haut nous
montre que $\underline{F}_{\bar{y}}$ \struck{est nul également} est nul
également si $\delta_{y} = 1$, par exemple si $\widetilde{X}$ admet une
section sur $\widetilde{Y}$. Nous allons supposer pour simplifier, par la
suite, que

\page{11}
\note{p.\ 19 du tapuscrit. Le premier tiers du feuillet, de (4.16) à
« des relations (4.15) », est encadré et barré de longs traits obliques ; il
porte lui-même des corrections de sa main, données ici}

\struck{(4.16) $H^{1}(y,\mathrm{Pic}_{X/k}) = 0$, $\mathrm{Br}(X_{y}) = 0$,
$H^{3}(y,\underline{\mathbb{G}}_{m}) = 0$} \add{$\mathrm{Ker}\bigl(H^{3}(y,
\mathbb{G}_{m}) \to H^{3}(X_{y},\mathbb{G}_{m})\bigr)$}
\note{le premier membre biffé est en partie une correction manuscrite,
lecture incertaine}

\struck{(Pour le voir, utiliser \add{l'homomorphisme} \struck{l'homomorphisme
de} suite\struck{s} exacte\struck{s} du type (4.5) relatives à $\bar{f} :
\bar{X} \to \bar{Y}$ \add{et à $f_{y} : X_{y} \to y$, qui s'écrit,}
\struck{\ill{}, qui donne,} compte tenu de (3.1) b) et de ( ) :}

\begin{tikzcd}[column sep=small, row sep=small, nodes={font=\scriptsize}]
\mathrm{Br}(\bar{Y}) \arrow[r] \arrow[d, "\alpha"] &
\mathrm{Br}(\bar{X}) \arrow[r] \arrow[d, "\beta"] &
H^{1}(\bar{Y},\bar{P}) \arrow[r] \arrow[d, "\gamma"] &
H^{3}(\bar{Y},\underline{\mathbb{G}}_{m}) \arrow[r] \arrow[d, "\alpha'"] &
H^{3}(\bar{X},\underline{\mathbb{G}}_{m}) \arrow[d, "\beta'"] \\
\mathrm{Br}(y) \arrow[r] & \mathrm{Br}(X_{y}) \arrow[r] &
H^{1}(y,P_{0}) \arrow[r] & H^{3}(y,\underline{\mathbb{G}}_{m}) \arrow[r] &
H^{3}(X_{y},\underline{\mathbb{G}}_{m})
\end{tikzcd}

\note{diagramme biffé avec le reste du bloc ; le carré de gauche est en
outre encadré à la main et barré}

\struck{Les flèches $\alpha$, $\beta$, $\alpha'$ du diagramme sont
bijectives, en vertu de (BR I ), (3.1), ( ), d'où résulte que $\gamma :
H^{1}(Y,P) \to H^{1}(y,P_{0})$ est injective, où $P_{0} =
\underline{\mathrm{Pic}}_{X_{y}/y}$. Donc (4.15) est une conséquence de
$H^{1}(y,P_{0}) = 0$, elle-même conséquence, d'après la deuxième suite exacte
ci-dessus, des relations (4.16).)} Alors l'image de $H^{1}(Y,P)$ dans
$H^{1}(Y,\underline{B})$ est le sous-groupe de $H^{1}(Y,\underline{B})$
formé des éléments dont l'image dans chaque $H^{1}(\bar{Y},\underline{B})$
est nulle. Ce sous-groupe s'appelle le \emph{groupe de Tate--Chafarévitch de
$Y$ à coéfficients dans $\underline{B}$}, et se note
$\mathrm{III}(Y,\underline{B})$ (cf plus bas \add{pour} la relation avec la
définition originelle de Tate et Chafarévitch). On trouve alors la suite
exacte déduite de (4.14) :
\[
\cdots\ \mathrm{Pic}(X_{\eta}/\eta) \to \coprod_{y}
H^{1}(y,\underline{E}_{y}) \longrightarrow H^{1}(Y,P) \longrightarrow
\mathrm{III}(Y,\underline{B}) \longrightarrow 0 ,
\tag{4.17}
\]
explicitant $H^{1}(Y,P)$ comme une extension du groupe de Tate--Chafarévitch
par un quotient de la somme $\coprod_{y} H^{1}(y,\underline{E}_{y})$.
\marginal{(P)}

Il faut enfin expliciter cette dernière somme. Notons que
$\underline{E}_{y}$ \struck{\ill{}} est nul si $X_{y}$ est
\struck{géométriquement} \add{géométriquement} intègre, à fortiori si $X_{y}$
est lisse sur $y$, de sorte que si $X_{\eta}$ est séparable sur $k(\eta)$
(par exemple, lisse sur $k(\eta)$), alors le support de $E$ est un
\struck{\ill{}} \add{contenu dans (en fait, égal à) l'ensemble} \emph{fini}
des $y_{i} \in Y$ tels que $X_{y_{i}}$ ne soit pas géométriquement intègre.
Un calcul facile donne pour $\underline{E}_{y}$ la structure suivante (en
procédant par descente à partir de $\underline{E}_{\bar{y}}$). On écrit
\add{l'égalité de \struck{\ill{}} diviseurs sur $X$} :
\[
X_{y} = \sum_{i} a_{y}^{i} C_{y}^{i} ,
\tag{4.18}
\]
les $C_{y}^{i}$ étant les composantes irréductibles de $X_{y}$, et \add{les}
$a_{y}^{i}$ leur multiplicité dans la fibre. On désigne par

\page{13}
\note{p.\ 20 du tapuscrit}

\[
\mu_{y}^{i} , \ \nu_{y}^{i}
\tag{4.19}
\]
la multiplicité radicielle et la multiplicité séparable de $C_{y}^{i}$ sur
$k(y)$ respectivement (EGA IV 4\add{.7.4.}), -- donc $\mu_{y}^{i} = 1$ si
$k(y)$ est parfait. On \struck{pose aussi} considère aussi la plus petite
extension galoisienne $k(y)^{i}$ \struck{(4.18) $Z_{y}^{i} =
\mathrm{Spec}(H^{0}(X_{y}^{i}, \mathcal{O}_{X_{y}^{i}}))$, $Z = \coprod
Z_{y}^{i}$ de sorte que l'on a simplement $Z_{y}^{i}$ est le spectre d'une
extension finie $k(y)_{i}$ de $k(y)$, de degré séparable} de $k(y)$ tel que
les composantes irréductibles de $(X_{y}^{i}) \otimes_{k(y)} k(y)^{i}$ soient
géométriquement irréductibles : c'est une extension finie \add{de $k(y)$,}
de degré $\nu_{y}^{i}$. On pose
\[
Z_{y}^{i} = \mathrm{Spec}(k(y)^{i}) , \qquad Z_{y} = \coprod Z_{y}^{i} .
\tag{4.20}
\]
Ceci posé, le faisceau $\underline{E}_{y}$ sur $y$ s'insère dans une suite
exacte de faisceaux
\[
0 \to \underline{\mathbb{Z}}_{y} \longrightarrow
p_{y*}(\underline{\mathbb{Z}}_{Z_{y}}) \longrightarrow \underline{E}_{y}
\longrightarrow 0 ,
\tag{4.21}
\]
où $p_{y} : Z_{y} \to y$ est la projection structurale, et où
l'homomorphisme $\underline{\mathbb{Z}}_{y} \to
p_{y*}(\underline{\mathbb{Z}}_{Z_{y}})$ est adjoint de l'homomorphisme
\[
p_{y}^{*}(\underline{\mathbb{Z}}_{y}) = \underline{\mathbb{Z}}_{Z}
\longrightarrow \underline{\mathbb{Z}}_{Z}
\tag{4.22}
\]
qui, sur la composante $Z_{y}^{i}$, se réduit à la multiplication par
$a_{y}^{i}$. La suite exacte de cohomologie associée à (4.21) nous fournit
alors
\begin{align}
H^{1}(y,\underline{E}_{y}) &= \mathrm{Ker}\bigl(H^{2}(y,\underline{\mathbb{Z}}_{y})
\longrightarrow \textstyle\prod H^{2}(Z_{y}^{i},\underline{\mathbb{Z}})\bigr)
\tag{4.23} \\
&= \mathrm{Ker}\bigl(H^{1}(y,\underline{\mathbb{Q}}/\underline{\mathbb{Z}})
\longrightarrow \textstyle\prod H^{1}(Z_{y}^{i},\underline{\mathbb{Q}}/
\underline{\mathbb{Z}})\bigr) , \notag
\end{align}
où la flèche $H^{1}(y,\underline{\mathbb{Q}}/\underline{\mathbb{Z}}) \to
H^{1}(Z_{y}^{i},\underline{\mathbb{Q}}/\underline{\mathbb{Z}})$ est
l'homomorphisme de transition évident, multiplié par $a_{y}^{i}$. Comme
$H^{1}(y,\underline{E}_{y})$ est l'intersection des noyaux de ces
homomorphismes, \struck{\ill{}} \add{et qu'un tel noyau} est manifestement
annulé par $a_{y}^{i}\nu_{y}^{i}$, on voit que $H^{1}(y,\underline{E}_{y})$
est annulé par
\[
d_{y} = \mathrm{pgcd}(a_{y}^{i}\nu_{y}^{i}) .
\tag{4.24}
\]
Lorsque $k(y)$ est fini, ou plus généralement pseudo-fini au sens de Serre,
i.e.\ son groupe fondamental $\mathrm{Gal}(\overline{k(y)}/k(y))$ isomorphe
à $\hat{\underline{\mathbb{Z}}}$, alors on trouve
\note{les numéros (4.19) à (4.24) sont corrigés à la main sur des numéros
tapés (4.18) à (4.23), et le numéro (4.18) du bloc biffé est l'ancien}

\section*{Cohomologie réelle de $B_{G}$ (page 15)}

\page{15}
$G$ groupe de Lie compact connexe
\[
\begin{aligned}
C_{(0)}^{p,q} &= \Omega^{q}(G^{p+1}/G) \\
&\ \ \cup \\
C_{(1)}^{pq} &= \omega^{(q)}(G^{p+1}/G) \\
&\ \ \downarrow \\
C_{(2)}^{pq} &= \omega_{\mathrm{inv}}^{(q)}(G^{p+1}/G)
\end{aligned}
\]
\marginal{\ill{} \uncertain{formes} \ill{} \uncertain{les dernières}
$E_{1}$ (\uncertain{cohomologie pour $V$ fixe})}
\note{à droite de la flèche verticale, « $i\delta$ », lecture incertaine}

Dans ce dernier bicomplexe, l'opérateur différentiel \uncertain{pour
variable} $q$ \emph{est nul}. (D\ill{})
\[
H^{*}(B_{G},\mathbb{R}) = H^{*}(C_{\bullet}^{pq}) \simeq H^{*}(C_{(1)}^{pq})
\simeq \sum_{p+q} \bigl\{ H^{p}(\nu \mapsto \omega_{\mathrm{inv}}^{(q)}
(G^{\nu+1}/G)) \bigr\}
\]
\[
\qquad = E_{2}^{pq}
\qquad \text{(isomorphismes canoniques)}
\]

\emph{La suite spectrale dégénère} : différentielles $d_{r}$ nulles pour
$r \geq 2$.

Soit $T \subset \omega_{\mathrm{inv}}^{*}(G)$, éléments transgressifs.
(\uncertain{OK} \ill{} des degrés impairs)

\ill{}.
\[
T \xrightarrow{\ \varphi\ } E_{2}^{1*} = \omega_{\mathrm{inv}}^{*}(G \times
G/G) \subset \omega_{\mathrm{inv}}^{*}(G \times G) =
\omega_{\mathrm{inv}}^{*}(G) \otimes \omega_{\mathrm{inv}}^{*}(G)
\]
où
\[
\varphi(T) = T \otimes 1 - 1 \otimes T
\]

\emph{Th}. $\mathrm{Sym}^{*}(T) \to H^{*}(B_{G},\mathbb{R})$ \emph{déduit
de $\varphi$} est \emph{un isomorphisme.}

\section*{Tapuscrit d'une lettre sur les catégories $C(X,A)$ (page 16)}

\page{16}
\note{feuillet dactylographié d'une lettre, sans en-tête ni signature sur ce
feuillet, barré de trois longs traits obliques ; une insertion de sa main.
Les soulignements tapés sous les lettres $\mathbb{Z}$ sont rendus par
$\underline{\mathbb{Z}}$}

triangulée $C(X,A)$ associée ainsi à $X$ détermine l'ensemble des
composantes connexes de $X$, et le groupe fondamental de chacune des
composantes connexes, et bien entendu la cohomologie de chacune, à
coéfficients des $A$-modules tordus de type fini \ldots{} Lorsque
$A = \underline{\mathbb{Z}}$, il semble bien que Hurewicz (ou une variante de
H.\ mise au point par Artin--Mazur) montre qu'un morphisme $f : X \to Y$ est
un isomorphisme dans la catégorie homotopique, si et seulement si le
foncteur $C(Y,\underline{\mathbb{Z}}) \to C(X,\underline{\mathbb{Z}})$ est
une équivalence. On se demande alors si on ne peut pas réconstituer
entièrement $X$ en termes de $C(X,\underline{\mathbb{Z}})$, du moins (si on a
fait dans la définition de $C(X,\underline{\mathbb{Z}})$ les hypothèses de
finitude que j'ai dites) si on se restreint à des $X$ satisfaisant certaines
conditions de finitude (comme des polyèdres qui sont finis en chaque
dimension). Plus précisément, serait-il vrai que tout foncteur exact
$C(Y,\underline{\mathbb{Z}}) \to C(X,\underline{\mathbb{Z}})$ proviendrait à
isomorphisme près d'un unique morphisme $X \to Y$ ? Même si ceci était trop
optimiste, les catégories $C(X,A)$ me semblent des invariants rupinants, qui
doivent être des approximations très bonnes de l'objet $X$ à homotopie près,
et on aimerait expliciter ce point de vue en \struck{expliciter}
interprétant en \struck{temmes} termes de ces catégories les constructions
les plus importantes familières aux homotopistes : construction des groupes
d'homotopie supérieurs, plus \struck{généralement} précisément des espaces
de lacets, et en sens inverse des espaces classifiants, ou des suspensions,
produits cartésiens d'espaces, \add{squelettes, images des groupes
d'homotopie supérieurs,} etc. \struck{Notes aussi que} C'est là un complexe
de questions manifestement très lié à celui que j'ai abordé dans ma lettre
du 26.3.\ dont tu viens de m'envoyer une photocopie, et je crois que
quelqu'un devrait bien un jour regarder ces choses systématiquement. Notes
que si tu fais sur $X$ une hypothèse de finitude comme plus haut, alors
$C(X,\underline{\mathbb{Z}})$ sera
\note{« rupinants » et « Notes que » sont tapés ainsi. Le mot biffé avant
« C'est là » est tapé puis barré à la machine ; la lecture « Notes aussi
que » est incertaine}

\section*{Groupes de Lie sans $p$-torsion (page 17)}

\page{17}
$G$ groupe de Lie \struck{compact} \add{\uncertain{complexe réductif}}
connexe. Considérons les conditions ($p \in \mathbb{P}$ donné)

\begin{itemize}
\item[(i)] $G$ sans $p$-torsion
\item[(ii)] $B_{G}$ sans $p$-torsion
\item[(iii)] $\mathrm{Sym}_{\mathbb{Z}}^{*}(\hat{B}) \to H^{*}(G/B,
\struck{\mathbb{Z}}\,\mathbb{Z})$ a une image d'indice premier à $p$, i.e.\
$\mathrm{Sym}_{\mathbb{F}_{p}}(\hat{B} \otimes_{\mathbb{Z}} \mathbb{F}_{p})
\to H^{*}(G/B,\mathbb{F}_{p})$ est surjectif.
\item[(i bis)] $H^{*}(G,\mathbb{F}_{p}) \simeq \bigwedge^{*} V$, $V$
contient \uncertain{générateurs} \add{de rang $r$} sur $\mathbb{F}_{p}$,
formé\add{s} d'éléments primitifs.
\item[(ii bis)] $H^{*}(B_{G},\mathbb{F}_{p})$ est une algèbre de polynômes
sur $\mathbb{F}_{p}$ à $r$ générateurs.
\end{itemize}

\struck{(i) N.B.\ Dans (ibis), on doit avoir} cf notes de Borel\ldots)

(i) $\Longrightarrow$ (i bis) par th.\ de Hopf.

(i bis) $\Longrightarrow$ (ii bis) par suite spectrale \struck{\ill{}}
\uncertain{descendante}, ou th.\ de transgression de Borel.

(ii bis) $\Longleftrightarrow$ (iii) [Suite spectrale de Leray dans
\uncertain{fibration} $G/B \to \mathbb{B}_{B} \to \mathbb{B}_{G}$ est
triviale] \struck{\ill{}}

(ii bis) $\Longrightarrow$ (ii) par égalité des nombres de Betti, coeff
$\mathbb{Q}$ ou $\mathbb{F}_{p}$

\begin{tikzcd}
(\mathrm{i}) \arrow[r, Rightarrow] \arrow[d, Rightarrow] &
(\mathrm{i\ bis}) \arrow[dl, Rightarrow] \\
(\mathrm{ii\ bis}) \arrow[r, leftrightarrow] \arrow[d, Rightarrow] &
(\mathrm{iii}) \\
(\mathrm{ii}) &
\end{tikzcd}

\note{la double flèche horizontale du bas est dessinée avec deux pointes ;
le trait vertical de (ii bis) à (ii) est double}

\section*{EGA IV, feuillet IV-1102 annoté (page 18)}

\page{18}
\note{feuillet du tapuscrit d'EGA IV, numéroté « IV-1102 » à la main.
Insertions tapées ou manuscrites données par \add{}, suppressions par
\struck{}, soulignements tapés par l'italique}

\add{(X étant supposé \emph{noethérien})}

(21.6.5) Considérons donc un diviseur \emph{positif} $D$ sur $X$, de sorte
que l'Idéal \add{fractionnaire} $I_{X}(D)$ est un \emph{Idéal de
$\mathcal{O}_{X}$}, qui est un $\mathcal{O}_{X}$-Module \emph{inversible} ;
soit $Y(D)$ le sous-préschéma fermé de $X$ qu'il définit. \struck{\ill{}}
\add{Pour} tout $x \in Y(D)$, \struck{\ill{}} il y a par hypothèse un
voisinage \add{ouvert} $U$ de $x$ dans $X$ et une section \emph{régulière}
$t \in \Gamma(U,\mathcal{O}_{X})$ telle que $I_{X}(D)|U =
\mathcal{O}_{U}.t$ (21.2.7) ; en d'autres termes, l'immersion canonique
$Y(D) \to X$ est \emph{régulière} \add{(19.1.4)} en tout point de $Y(D)$
\emph{et de codimension 1}. On notera d'ailleurs qu'inversement, si $Y$ est
un sous-préschéma fermé de $X$, régulièrement immergé dans $X$ et
\struck{\ill{}} de codimension 1 en tout point de $Y$, il existe un diviseur
positif et un seul \struck{\ill{}} $D$ sur $X$ tel que $Y(D) = Y$, car
\add{pour} \struck{\ill{}} tout $x \in Y$, il y a un voisinage $U$ de $x$
dans $X$ tel que $Y \cap U$ soit défini par un Idéal de $\mathcal{O}_{U}$ de
la forme $\mathcal{O}_{U}.t$, où $t$ est \emph{régulier} dans
$\Gamma(U,\mathcal{O}_{X})$.
\marginal{devrait figurer dans 21.2. \ill{} (sans \uncertain{noeth.})}
\note{la note marginale est reliée par un trait au passage « On notera
d'ailleurs \ldots{} $\Gamma(U,\mathcal{O}_{X})$ »}

Puisque $X$ est supposé noethérien, $Y(D)$ n'a qu'un nombre fini de points
maximaux, et en chacun de ces points $x$, on a $\mathrm{codim}(\overline{\{x\}},X)
= \dim(\mathcal{O}_{X,x}) = 1$, autrement dit $x \in X^{(1)}$ ;
\struck{\ill{}} en outre, $\mathcal{O}_{Y(D),x}$ est un anneau
\emph{artinien} en chacun de ces points. En tout point $x \in X^{(1)}$ qui
n'est pas point maximal de $Y(D)$, on a nécessairement $x \notin Y(D)$, donc
$\mathcal{O}_{Y(D),x} = 0$. Posons
\marginal{réf !}
\[
\mathrm{cyc}(D) = \sum_{x \in X^{(1)}} \mathrm{long}(\mathcal{O}_{Y(D),x})
\cdot \overline{\{x\}}
\tag{21.6.5.1}
\]
qui est donc un élément de $Z^{1}(X)$.

\emph{Proposition} (21.6.6). -- \struck{Si $X$ est noethérien,}
\emph{l'application $D \mapsto \mathrm{cyc}(D)$ est un homomorphisme du
monoïde $\mathrm{Div}^{+}(X)$ dans $Z^{1}(X)$.}

Tout revient à voir que pour deux diviseurs \emph{positifs} $D$, $D'$, on a
$\mathrm{cyc}(D+D') = \mathrm{cyc}(D) + \mathrm{cyc}(D')$. Or, on a (21.2.5)
$I_{X}(D+D') = I_{X}(D) I_{X}(D')$ ; tout revient à voir que si $x \in
X^{(1)}$, \add{si l'on pose} $A = \mathcal{O}_{X,x}$, et si $t$ et $t'$ sont
deux éléments réguliers de $A$, on a $\mathrm{long}(A/tt'A) =
\mathrm{long}(A/tA) + \mathrm{long}(A/t'A)$ ; mais puisque \struck{\ill{}}
$t$ est régulier, $tA/tt'A$ est isomorphe à $A/t'A$, d'où \struck{\ill{}}
aussitôt la proposition.

Il résulte aussitôt des définitions que pour tout ouvert $U \subset X$, on

\end{document}
